\chapter{分子数较少时：从浓度轨迹转向概率分布}
前面的轨迹把表达量当作连续浓度。在分子数较多、涨落相对较小时，这能突出反馈的主要作用；但同一个产生--降解机制在低分子数下也会不断随机变化。此时，长期“稳定”不再意味着每条轨迹停在同一个数值，而是分子数的概率分布保持不变。先回到最简单的无反馈反应，看看这一转换怎样进行。

\section{主方程、稳态分布与涨落的时间结构}
考虑 $k>0,\gamma>0$ 的出生--死亡过程
\begin{equation}\varnothing\xrightarrow{k}X,\qquad X\xrightarrow{\gamma}\varnothing.\end{equation}
令 $P_n(t)$ 表示时刻 $t$ 恰有 $n$ 个分子的概率。在短时间 $\dd t$ 内，出生使 $n-1\to n$，跃迁率为 $k$；死亡使 $n+1\to n$，跃迁率为 $\gamma(n+1)$；从 $n$ 离开的总率为 $k+\gamma n$。流入减去流出，只保留 $\dd t$ 的一阶项，就得到
\begin{equation}\dot P_n=kP_{n-1}+\gamma(n+1)P_{n+1}-(k+\gamma n)P_n,\qquad P_{-1}=0.\end{equation}
零边界没有向负分子数的跃迁。确定性方程描述一个浓度的导数，而主方程同时描述所有可能分子数之间的概率转移。


先寻找长期不变的分布。定义从 $n$ 到 $n+1$ 的净概率流
$j_n=kP_n-\gamma(n+1)P_{n+1}$，主方程就是 $\dot P_n=j_{n-1}-j_n$。在稳态各相邻概率流必须相等，而零边界给出 $j_{-1}=0$，所以所有 $j_n=0$。这就推出相邻态的\emph{详细平衡}：每一对相邻分子数之间，向上与向下转移的概率流分别相等。它比只要求总分布不再变化更具体。递推得到
\[
P_n=P_0\frac{\mu^n}{n!},\qquad\mu=\frac k\gamma.
\]
归一化要求 $1=P_0\sum_{n=0}^{\infty}\mu^n/n!=P_0e^\mu$，从而
\begin{equation}P_n=e^{-\mu}\frac{\mu^n}{n!}.\end{equation}
因此泊松分布是这条链的唯一归一化稳态；平均分子数 $\mu$ 与此前产生--降解的平衡值对应，但单次观测仍可偏离平均值。
\begin{figure}[H]\centering\includegraphics[width=.74\linewidth]{figures/birth_death.pdf}\caption{基本出生--死亡模型的稳态泊松分布。平均分子数越低，涨落相对于平均值越明显。}\end{figure}


稳态存在还不等于已经说明任意初始分布会到达它。在线性出生--死亡模型中，这一步可以精确补上：从确定的 $n_0$ 个分子出发，每个初始分子存活到 $t$ 的概率为 $p_t=e^{-\gamma t}$，所以存活数服从 $\operatorname{Bin}(n_0,p_t)$。时刻 $s$ 出生的分子存活到 $t$ 的概率为 $e^{-\gamma(t-s)}$，把泊松出生流按这个存活概率保留，新生且仍存活的分子数服从参数为
\[
\int_0^t ke^{-\gamma(t-s)}\dd s=\mu(1-e^{-\gamma t})
\]
的泊松分布，与初始存活数独立。因此
\[
n(t)\overset{\rm law}=\operatorname{Bin}(n_0,e^{-\gamma t})
 +\operatorname{Pois}\!\left(\mu(1-e^{-\gamma t})\right).
\]
记号表示两边具有相同的概率分布。随时间增长，初始分子最终消失，新生存活数趋于 $\operatorname{Pois}(\mu)$。对任意几乎必然有限的随机初始分子数，先条件于 $n_0$ 再平均也得到相同极限。这证明该泊松分布不仅唯一，而且吸引所有这类初始分布。

为了比较主方程与浓度模型，可以直接计算平均值与方差。先假设平均值与方差都是有限的，数学上称为\emph{有限二阶矩}：$\langle n^2\rangle<\infty$，所以分子数平方的平均值不会无限大。下面分别关心分子数 $h(n)=n$ 和它的平方 $h(n)=n^2$。把每次出生、死亡对这些数量的改变，按照发生概率加起来，就有
\[
\frac{\dd\langle h(n)\rangle}{\dd t}
=\left\langle k[h(n+1)-h(n)]+\gamma n[h(n-1)-h(n)]\right\rangle.
\]
先取 $h=n$，得到
\begin{equation}\frac{\dd\langle n\rangle}{\dd t}=k-\gamma\langle n\rangle.\end{equation}
这个平均值方程与确定性产生--降解模型具有相同形式，是线性反应下的精确关系。再取 $h=n^2$，得到
\[
\frac{\dd\langle n^2\rangle}{\dd t}=2k\langle n\rangle+k-2\gamma\langle n^2\rangle+\gamma\langle n\rangle.
\]
减去 $\dd\langle n\rangle^2/\dd t$，就有
\begin{equation}\frac{\dd\operatorname{Var}(n)}{\dd t}=k+\gamma\langle n\rangle-2\gamma\operatorname{Var}(n).\end{equation}
在稳态代入 $\langle n\rangle=\mu=k/\gamma$，得到 $\operatorname{Var}(n)=\mu$。因此变异系数为 $1/\sqrt\mu$，分子数越低，相对内禀涨落越大。

涨落还具有时间结构。由平均值方程，从给定 $n(0)$ 出发的条件均值为
\[
\mathbb E[n(t)\mid n(0)]=\mu+(n(0)-\mu)e^{-\gamma t}.
\]
减去 $\mu$，再乘以 $n(0)-\mu$ 并在稳态平均，得到
\begin{equation}\operatorname{Cov}(n(t),n(0))=\mu e^{-\gamma t},\qquad t\geq0.\end{equation}
所以同一个 $1/\gamma$ 既描述平均浓度的恢复，也描述这条线性反应链的涨落相关时间。若产生率带有非线性反馈，通常有 $\langle F(n)\rangle\ne F(\langle n\rangle)$，平均值不再自动满足原来的非线性浓度方程。

\section{随机切换、双峰与记忆的寿命}
前面的自激活和互抑制开关有两个确定性吸引态，但在随机模型中，两个概率峰还可能来自另一种时间结构：启动子在两种产生率之间缓慢切换。取关闭与开启状态的产生率为 $k_0,k_1$，共享降解率 $\gamma$，切换率为 $k_{\rm on},k_{\rm off}$。

如果两种切换均远慢于 $\gamma$，蛋白质常有足够时间接近当前启动子状态的条件平衡。二态链的稳态流量平衡给出开启概率 $q=k_{\rm on}/(k_{\rm on}+k_{\rm off})$，整体分布的领先阶近似于是为
\[
P_n\simeq(1-q)\operatorname{Pois}_{k_0/\gamma}(n)+q\operatorname{Pois}_{k_1/\gamma}(n).
\]
若两个均值相距足够远且两权重均显著，这个混合分布可以双峰；每个固定启动子状态下的浓度模型却都只有唯一稳定的稳态。该式是慢切换近似，有限切换速率下不是精确公式。由此可见，解释两个峰时还需要看时间序列、驻留时间和扰动响应，判断状态选择来自确定性反馈，还是来自慢切换或其他机制。

真正的确定性双稳态进入随机模型后，记忆通常具有有限寿命。若出生率依赖当前分子数，写成 $\lambda_n>0$，死亡率仍为 $\gamma n$，详细平衡递推变成
\[
\pi_n=\pi_0\prod_{j=0}^{n-1}\frac{\lambda_j}{\gamma(j+1)}.
\]
当出生率有上限时，这个乘积由泊松尾部控制，可以归一化；因为所有相邻跃迁均可发生，整条链仍只有一个归一化稳态。两个概率峰可以对应两个长时间驻留区，但不代表两份互不相通的稳态分布。连续模型的吸引域分界在这里可以被一连串随机出生或死亡跨过，所以“双稳态记忆”变成很长但有限的驻留时间。

\noindent\begin{minipage}{\linewidth}\setlength{\parindent}{2em}
这个寿命可以用\emph{首次到达时间}具体定义。选择两个分子数 $i<j$，将 $j$ 作为切换判据，令 $T_i$ 为从 $i$ 出发第一次到达 $j$ 所需的平均时间。对 $k<j$，按第一次跃迁前的平均等待时间分解，有
\[
-1=\lambda_k(T_{k+1}-T_k)+\gamma k(T_{k-1}-T_k),\qquad T_j=0.
\]
令 $\Delta_k=T_k-T_{k+1}$，并用 $\lambda_{k-1}\pi_{k-1}=\gamma k\pi_k$，上式乘以 $\pi_k$ 后可逐项相加：
\[
\lambda_k\pi_k\Delta_k=\sum_{m=0}^{k}\pi_m.
\]
因此得到精确的离散表达式
\[
T_i=\sum_{k=i}^{j-1}\frac{\sum_{m=0}^{k}\pi_m}{\lambda_k\pi_k}.
\]
若两个峰之间的概率很小，分母中的 $\pi_k$ 会让相应等待时间变大。可见解释记忆时，除了是否有两个稳定表达态，还要问噪声有多强、跨越阈值需要多久，以及这个时间是否长于实际观察窗口。\end{minipage}\par
这里的公式适用于一维出生--死亡链；双基因联合随机过程需要在二维分子数空间重新求首次到达问题，不能把这条一维公式直接照搬过去。
